\subsection{An elementary isogeny descent for Type 16}\label{detail:type16}
We give the complete descent behind the Type~16 exclusion, including its
exceptional charts. In this type $c$ is the squared singleton center on
$B$, and $S,T$ and $U,V$ are respectively the sum and product of the
singleton squares on $A$ and the doubleton squares on $B$. The valid
moment branch is
\[
 S=\frac{195-20c}{12},\quad T=\frac{160c^2-1335c+2304}{36},
 \quad U=\frac{4(9-c)}3,\quad V=\frac{2(c-6)(4c-9)}3.
\]
Here are the divisions used to reach that branch. The first two moments give
\[
 U=\frac{4S-4c+31}{8},\quad
 V=\frac{-16S^2-32Sc+296S+64T+48c^2-296c-73}{128}.
\]
Writing $D=4S+4c-51$, the third moment determines $T$ when $D\ne0$.
The fourth residual then factors as a nonzero constant times
\[
 (4c-25)(4c-1)\prod_{j\in\{1,9,49\}}(-4S+4c+j)
 (12S+20c-195)/D^2.
\]
At $D=0$ the third residual is $-3(4c-25)(4c-21)(4c-1)/64$.
The values $c=1/4,25/4$ are central overlaps, and $21/4$ is nonsquare.
For $S=c+j/4$ the $A$ singleton polynomial is $(X-c)(X-j/4)$,
so the two colors share the singleton center. Thus every valid input
belongs to the displayed branch.

Since $c=v^2$ and $V$ is square, it gives a rational point on
\[
 \mathcal B:\quad y^2=\frac23(v^2-6)(4v^2-9).
\]
We claim that its complete projective rational-point set is
\[
 \mathcal B(\Q)=\{(3/2,0),(-3/2,0),(0,6),(0,-6)\}.
\]
There are no rational infinities, because the leading coefficient $8/3$
is nonsquare. With $t=v-3/2$ the equation becomes
$y^2=-30t+14t^2+16t^3+(8/3)t^4$. The substitutions
\[
 X=-10-30/t,\quad Y=-30y/t^2,
 \qquad v=3/2-\frac{30}{X+10},\quad y=-\frac{30Y}{(X+10)^2}
\]
give the elliptic curve
\[
 E:Y^2=X^3+44X^2+100X,
 \qquad E':Y^2=X^3-88X^2+1536X.
\]
The omitted point $(3/2,0)$ corresponds to the identity. A point of $E$
with $X=-10$ would require $Y^2=2400$, which is impossible.

For $E_{A,B}:Y^2=X(X^2+AX+B)$ the two-isogeny descent map has values
\[
 \alpha(O)=1,\qquad \alpha((0,0))=B,\qquad
 \alpha((X,Y))=X\quad\text{in }\Q^\times/\Q^{\times2}.
\]
Its kernel is the image of the dual isogeny. Taking valuations shows that
each image has a signed squarefree representative $d\mid B$. A point
with $X=du^2/v^2$ in lowest terms gives
\begin{equation}\label{eq:early-isogeny-quartic}
 w^2=du^4+Au^2v^2+(B/d)v^4,\qquad u,v,w\in\Z,
 \quad\gcd(u,v)=1.
\end{equation}
This follows by putting $Y=duw/v^3$ in the elliptic equation.
The square of a rational number is integral only if that number is
integral. The isogeny exact sequences and the Mordell--Weil theorem give
\begin{equation}\label{eq:early-isogeny-rank}
 2^{\operatorname{rank}E(\Q)}
 =\frac{|\alpha(E(\Q))|\,|\alpha'(E'(\Q))|}{4}.
\end{equation}
Both isogeny kernels have order two. This is the standard two-isogeny
rank formula~\cite[\S2]{FisherHigher}.

For $A=44,B=100$, the classes are $\pm1,\pm2,\pm5,\pm10$.
The following are the complete right-side residue sets in the primitive
charts for \eqref{eq:early-isogeny-quartic}.
\begin{center}
\begin{tabular}{ccl}
\toprule
$d$ & Modulus & Residues\\
\midrule
$-10,-1$ & $9$ & $6,8$\\
$2,5$ & $9$ & $2,5,6$\\
$-5$ & $16$ & $3,11,12$\\
$-2$ & $16$ & $8,14$\\
\bottomrule
\end{tabular}
\end{center}
For $p^e$ the charts are $(u,1)$ for all residues $u$, and $(1,v)$
with $p\mid v$. They exhaust primitive pairs after a unit rescaling.
Such a rescaling multiplies the right side by a fourth power, preserving
squarehood. The table avoids the square sets $\{0,1,4,7\}$ modulo 9 and
$\{0,1,4,9\}$ modulo 16. The remaining classes 1 and 10 occur at
$O$ and $(10,80)$, so $\alpha(E(\Q))=\{1,10\}$.
On $E'$ the classes are $\pm1,\pm2,\pm3,\pm6$. If $d<0$, all three
coefficients in \eqref{eq:early-isogeny-quartic} are negative, so there
is no nonzero real solution. For $d=2,3$ the complete primitive
residue set modulo 5 is $\{2,3\}$. The remaining classes 1 and 6
occur at the isogeny-kernel points, since $1536=6\cdot16^2$.
Formula~\eqref{eq:early-isogeny-rank} proves rank zero.

The discriminant of $E$ is $2^{17}\cdot3\cdot5^4$. Direct counts give
$\#E(\F_7)=4$ and $\#E(\F_{11})=12$. Torsion of
order prime to the reduction characteristic injects at good reduction.
Using both primes also removes the exceptional 7- and 11-primary
possibilities, so the full rational torsion order divides four.
The four points $O,(0,0),(10,\pm80)$ attain this bound, with
$2(10,80)=(0,0)$. Their inverse images are exactly the four stated
points of $\mathcal B$.

At $c=0$ one has $S=65/4,T=64$, and the singleton discriminant is
$S^2-4T=129/16$, a nonsquare. At $c=9/4$ the singleton squares are
$49/4,1/4$; the latter repeats vertices of the central chain of length
six. Hence no rational point gives a valid Type~16 input.
\cert{check_type16_exclusion.py} verifies the maps, residue tables,
reduction counts, doubling identity and final substitutions.

\subsection{Scalar normalization and the two elliptic images for Type 5}\label{detail:type5}
Here the valid moment branch is
\[
 S=\frac{20c-85}{16},\quad T=\frac{80c^2-840c+589}{256},
 \quad U=\frac{12c-83}{16},\quad V=\frac{(4c-49)(4c-9)}{256}.
\]
The factor $c=1/4$ gives a central overlap. At $c=9/4$ the $A$
paired-center polynomial has root 1 and repeats central vertices.
At $c=25/4$ it has root 4, making an $A$ vertex equal to the $B$
singleton. These are all constant factors removed when determining
$S,T,U,V$.
Retain the separate center conditions, which say that $c$ and
\[
 q_1(c)=80c^2-840c+589,\quad q_2(c)=80c^2-40c+4869,
 \quad q_3(c)=80c^2-1064c+5125
\]
are rational squares. The three quadratics retain the product and
both paired-center discriminants separately. Their product defines a
smooth genus-two curve. Put $f=q_1q_2q_3$ and
$\mathcal A=\Q[\theta]/(f(\theta))$.
The exact fake two-cover descent has two surviving classes represented
by $c_0-\theta$ for $c_0=49/4,81/4$.

We first show that a genuine input in either class satisfies
\begin{equation}\label{eq:type5-normalized-class}
 c-\theta=(c_0-\theta)z^2\quad\text{in }\mathcal A,
\end{equation}
without a rational nonsquare multiplier. Suppose initially that the
right side is $\lambda(c_0-\theta)z^2$, where
$\lambda\in\Q^\times$. Nonnegativity of $U,V$ gives
$c\geq83/12$ and $(4c-49)(4c-9)\geq0$, hence $c\geq49/4$.
The values of $(q_1,q_2,q_3)$ at the two anchors are
\[
 (48^2,128^2,64^2),\qquad(128^2,192^2,128^2).
\]
Each quadratic is irreducible, so all root differences used below are
nonzero at rational $c$ and $c_0$.

For an odd $p\ne5$, at least one of the three discriminants
\[
 2^{10}\cdot505,\qquad-2^{14}\cdot95,\qquad-2^{14}\cdot31
\]
is a $p$-unit. The corresponding quadratic, with leading coefficient
80, defines an unramified quadratic algebra at $p$. Write $r$ for
its root. If $v_p(c)<0$, the valuation is even because $c$ is square,
and every valuation of $c-r$ equals it. If $c$ is integral and the
factor is inert, $c-r$ is a unit. If it splits, its two roots have
distinct reductions, so at most one root difference has positive
valuation. Their product is $q_i(c)/80$, of even valuation, so both
valuations are even. Apply the same reasoning to $c_0$.
The squareclass identity now forces $v_p(\lambda)$ even.

At 5, a negative $v_5(c)$ would give $80c^2$ uniquely least valuation
in $q_1(c)$, namely $1+2v_5(c)$, which is odd. Thus $c$ is integral
at 5. The two Newton segments of $q_3$ give roots of valuations 3
and $-1$ in $\Q_5$. At the latter root,
\[
 \frac{c-r}{c_0-r}=\frac{1-c/r}{1-c_0/r}\in1+5\Z_5.
\]
This is a square, and the identity makes $\lambda$ a local square.

At 2 one must have $v_2(c)=-2$. A nonnegative valuation makes
$q_1(c)\equiv5\pmod8$. If the even valuation were at most $-4$,
the ratio $q_1(c)/(80c^2)$ would be 1 modulo 8: its two other terms
have relative valuations at least 3 and 4. The normalized leading
unit is again 5 modulo 8, a contradiction. Since 505 is a dyadic
square, $q_1$ splits with roots
\[
 r_\pm=\frac{105\pm4\sqrt{505}}{20},\quad
 v_2(r_\pm)=-2,\quad v_2(r_+-r_-)=1,\quad4r_\pm\equiv1\pmod8.
\]
For the square $c$, $4c\equiv1\pmod8$. Both $c-r_\pm$ have
valuation at least one, and at least one has valuation exactly one.
Their product $q_1(c)/80$ has even valuation, so both valuations
are odd. The corresponding assertion for $c_0$ makes both root
ratios have even valuation and forces $v_2(\lambda)$ even.
Finally both real roots of $q_1$ are below $49/4$, since
$\sqrt{505}<35$. The ratios are positive at both embeddings, so
$\lambda>0$. A positive rational with even valuation at every
prime is square. Absorbing its square root proves
\eqref{eq:type5-normalized-class}.

The normalization uses individual original squares. Projection to an
odd-degree factor before normalization would retain the multiplier
$\lambda$ to an odd power and would not define a fixed elliptic twist.
After normalization the following fixed twists are justified.
For $c_0=49/4$, take $K=\Q(r)$, $q_1(r)=0$, and
\[
 g(c)=\frac{(c-r)q_3(c)}{80},\qquad
 \gamma=g(49/4)=\frac{256}{5}(49/4-r).
\]
The norm of \eqref{eq:type5-normalized-class} on the cubic factor
gives $g(c)/\gamma$ square in $K$, hence $z^2=\gamma g(c)$.
If $g=c^3+g_2c^2+g_1c+g_0$, the monic model and parameter map are
\[
 E:v^2=u^3+\gamma g_2u^2+\gamma^2g_1u+\gamma^3g_0,
 \qquad (u,v)=(\gamma c,\gamma z),\quad \phi(u,v)=(u:\gamma).
\]
The point $P=(49\gamma/4,\gamma^2)$ has infinite order.
The arithmetic certificate gives rank one over $K=\Q(\sqrt{505})$
and full torsion order two. Two-saturate $P$ modulo torsion, then adjoin
the entire torsion subgroup. The resulting subgroup $G$ has finite odd
index in $E(K)$. Elliptic Chabauty with input index bound 2 returns
three elements, with image set $\{49/4,\infty\}$ and index obstruction
16. The actual odd index is coprime to 16, so these images are complete
for $E(K)$. The original parameter is finite, and the sextic also has
no rational infinity because its leading coefficient 512000 is nonsquare.
At $c=49/4$ one has $U=4,V=0$, so a reflected doubleton pair repeats.

For $c_0=81/4$, take $q_2(r)=0$, $K=\Q(\sqrt{-95})$ and
$\gamma=c_0-r$. Put $h(c)=c(c-r)q_3(c)$. The normalized cubic
norm and the additional square $c/c_0$ show $h(c)/h(c_0)$ square.
Since $q_3(c_0)=128^2$ and $c_0=(9/2)^2$, $h(c_0)$ has class
$\gamma$. Thus $z^2=\gamma h(c)$ and its seed ordinate is
$z_0=\gamma(9/2)128$. With $b=-1064,d=5125,L=-\gamma rd$,
inversion followed by scaling gives
\[
 E:v^2=u^3+\gamma(d-rb)u^2+L\gamma(b-80r)u+80L^2\gamma,
 \qquad u=L/c,\quad v=Lz/c^2,\quad\phi(u,v)=(L:u).
\]
The real bound $c\geq49/4$ ensures the inversion loses no input.
The seed $(L/c_0,Lz_0/c_0^2)$ is non-torsion, the rank is one,
and torsion has order two. The same construction of a finite odd-index
subgroup gives three Chabauty elements, obstruction 8, and complete
images $\{0,81/4\}$. At zero $S=-85/16$ is impossible. At $81/4$
the paired-center squares are $16,4$ on $A$ and $9,1$ on $B$,
with singleton centers $\pm9/2$ on $B$. Both inputs are the same ten
half-integers. This completes both classes.
All class-group proof maps are actual Minkowski bounds, including
fields introduced inside the elliptic calculations.
\cert{check_type5_exclusion.m} carries out the descents, rank and
index computations. \cert{check_type5_scalar_normalization.py}
checks the polynomial, local root and Newton identities.

\subsection{The ideal and unit calculation for Type 15}\label{detail:type15}
For Type~15 the noncollision moment formulas are
\[
 S=\frac{20c+15}{8},\quad T=\frac{160c^2-20c+49}{96},
 \quad U=\frac{69-4c}{8},\quad V=\frac{(2c-23)(4c-31)}{48}.
\]
The removed values $c=1/4,9/4,25/4$ meet opposite-color chains.
The rational quantities $c$ and
\[
\begin{aligned}
 q_T&=6(160c^2-20c+49),&q_A&=3(-80c^2+1960c+283),\\
 q_B&=3(-80c^2+808c+2875),&q_V&=3(2c-23)(4c-31)
\end{aligned}
\]
are squares. The latter are respectively $576T$, $576(S^2-4T)$,
$576(U^2-4V)$ and $144V$.
Put
\[
 K=\Q(\sqrt{15}),\quad r=\frac{245+64\sqrt{15}}{20},
 \quad c_0=1/4,\quad R=\frac{c-r}{c_0-r}.
\]
The polynomial $q_A$ is irreducible. Its positivity places $c$ between
its real roots, as is $c_0$ because $q_A(c_0)=48^2$.
Thus $R$ is totally positive, and
$N_{K/\Q}(R)=q_A(c)/q_A(c_0)$ is a rational square.

Outside 2,3,5 the polynomial $q_A$ has unit leading coefficient and
discriminant. At a pole of $c$ its even valuation gives an even
valuation of both root differences. For integral $c$ an inert root
difference is a unit; in the split case at most one root difference
has positive valuation, which must be even because the rational norm
is square. The same holds for $c_0$. At 2,3,5 the field has one
prime of residue degree one, so the square rational norm directly
forces its valuation of $R$ to be even. Consequently $(R)=I^2$.

The ring of integers is $\Z[\sqrt{15}]$ and its discriminant
is 60. Its Minkowski bound is $\sqrt{15}<4$, so the class group
is generated by the unique primes $\mathfrak p_2,\mathfrak p_3$
of norms two and three. Since
$N(3+\sqrt{15})=-6$, one has $(3+\sqrt{15})=\mathfrak p_2\mathfrak p_3$,
and $\mathfrak p_2^2=(2)$. The prime $\mathfrak p_2$ is
nonprincipal: norm $\pm2$ would require a rational integer square
to be 2 or 3 modulo 5. The class group therefore has order two.
It follows that
\[
 R=2^d\epsilon^e z^2,\qquad d\in\{0,1\},\quad
 \epsilon=4+\sqrt{15},\quad z\in K^\times.
\]
Here the totally positive units are exactly the powers of $\epsilon$.
There is no norm $-1$ unit by reduction modulo 3. If a norm-one unit
$v=x+y\sqrt{15}$ satisfied $1<v<\epsilon$, then
$y=(v-v^{-1})/(2\sqrt{15})$ would be an integer strictly between zero
and one. Thus $\epsilon$ is fundamental.

At 5, a negative even valuation of $c$ would make the leading term
of $q_T(c)$ have the unique least valuation $1+2v_5(c)$, an odd
number. Hence $c$ is integral. The Newton polygon of $q_A$ has one
segment of length two and its roots have valuation $-1/2$ at 5.
Therefore $R\in1+\mathfrak p_5$, which is square in the completion.
The residue of $\epsilon$ is 4 in $\F_5$, whereas 2 is
nonsquare, forcing $d=0$. Finally
\[
 (3+\sqrt{15})^2=6\epsilon,\qquad
 \frac{49/4-r}{6(1/4-r)}=
 \left(\frac{2}{3+\sqrt{15}}\right)^2.
\]
It follows that either $(c-r)/(1/4-r)$ or $(c-r)/(49/4-r)$ is square
in $K$. This is a complete squareclass conclusion, with no height
bound on $c$.

For the two classes use respectively $(c_0,q,s_0)=(1/4,q_B,96)$ and
$(49/4,q_V,9)$, so $q(c_0)=s_0^2$. If $\gamma=c_0-r$, the original
squares give $z^2=\gamma(c-r)q(c)$. Write $q(c)=ac^2+bc+d$ and
$L=a\gamma$. Then
\[
 E:v^2=u^3+\gamma(b-ar)u^2+L\gamma(d-br)u-L^2\gamma dr,
 \qquad (u,v)=(Lc,Lz),\qquad\phi(u,v)=(u:L).
\]
The seed $(Lc_0,L\gamma s_0)$ has infinite order. Each curve has
rank one over $K$. Two-saturation of the seed and addition of the full
torsion subgroup, of orders two and four, gives a subgroup of finite
odd index. Elliptic Chabauty with input index bound 2 returns
\begin{center}
\begin{tabular}{cccc}
\toprule
$c_0$ & Number of group elements & Index obstruction & Rational images\\
\midrule
$1/4$ & $3$ & $4$ & $1/4,\infty$\\
$49/4$ & $5$ & $4$ & $31/4,23/2,49/4,\infty$\\
\bottomrule
\end{tabular}
\end{center}
The odd subgroup indices are coprime to the returned obstructions, so the
lists are complete. Every class-group proof bound in these computations
is a Minkowski bound. The images $31/4,23/2$ are nonsquares, and a
center parameter is finite. At $1/4$ an $A$ singleton is $1/2$,
already on its central chain. At $49/4$ the $A$ singleton absolute
centers are $9/2,7/2$, the $B$ singleton absolute centers are $3/2,1/2$,
and the $B$ tripleton centers are $\pm7/2$. Both inputs are the ten
half-integers from $-9/2$ to $9/2$. This excludes the two classes.
The complete arithmetic and exact unit identities are
\cert{check_type15_exclusion.m} and
\cert{check_type15_unit_normalization.py}.

\subsection{Component norms and the scalar for Type 4}\label{detail:type4}
For Type~4 the valid moment branch has
\[
 S=\frac{20c+25}{6},\quad T=\frac{240c^2+40c+167}{48},
 \quad U=\frac{44c+91}{12},\quad V=\frac{256c^2+97c+466}{36}.
\]
To check the exceptional factors directly, the first two moments give
\[
 U=2S-3c-3/4,\quad
 V=(8S^2-48Sc-24S+16T+48c^2+66c+9)/8.
\]
The resultant of the remaining two equations in $T$ is proportional to
\[
\begin{split}
 &(4c-9)(4c-1)(-6S+20c+25)(-2S+4c+1)\\
 &\qquad\cdot(4S^2-16Sc-20S+16c^2+24c+25).
\end{split}
\]
The first two factors are central overlaps. On $S=2c+1/2$, the third
moment gives $T=(4c-1)^2/16$. With $c=v^2$ the paired centers are
$\pm(v-1/2),\pm(v+1/2)$, whose doubletons overlap the $B$ tripleton.
On the quadratic branch,
\[
 S=2v^2+5/2\pm2v,\qquad
 T=(2v\mp1)^2(2v\pm3)^2/16,
\]
and the same overlap occurs. The coefficients used to determine $T$
can vanish only at $c=-11/4,1/4,9/4$; the negative value is not a
rational square and the other two were excluded. Thus the stated
branch loses no valid input.
The separate square conditions are $c$ and
\[
 q_T=3(240c^2+40c+167),\qquad
 q_A=-80c^2+220c+31,\qquad
 q_B=3(-720c^2+2152c+275).
\]
These are $144T$, $9(S^2-4T)$ and $144(U^2-4V)$.
Let $f=q_Tq_Aq_B$, $\mathcal A=\Q[\theta]/(f(\theta))$,
and $c_0=1/4$. The values of the three quadratics at $c_0$ are
$24^2,9^2,48^2$.

The fake two-cover descent of $Y^2=f(c)$ returns three classes.
For a representative $\delta$, let $\delta_i$ be its projection
to the quadratic factor field $K_i$ of $q_i$. Necessarily
\[
 \operatorname{lc}(q_i)N_{K_i/\Q}(\delta_i)
       \in\Q^{\times2}\qquad(i=T,A,B).
\]
The condition is independent of the fake representative, since a
rational scalar has square norm on a quadratic component.
For an actual point it is exactly the original square condition
$q_i(c)\in\Q^{\times2}$. Exactly one computed class passes
all three tests and it is represented by $c_0-\theta$.
Thus a genuine input has
$c-\theta=\lambda(c_0-\theta)z^2$ in $\mathcal A$.

The three discriminants are
\[
 -2^{10}3^2\cdot5\cdot31,\quad 2^4 3^6\cdot5,
 \quad 2^{14}3^2\cdot331.
\]
For each $p\notin\{2,3,5\}$ choose a component with unit leading
coefficient and discriminant. The same root-difference argument used
in Type~5 gives even valuations of both corresponding ratios. It
forces $v_p(\lambda)$ even.
If $v_3(c)<0$, it is even and at most $-2$. In
$q_T(c)=720c^2+120c+501$ the leading term is uniquely minimal,
of valuation $2+2v_3(c)$ and normalized unit $80\equiv2\pmod3$.
The other terms have relative valuations at least one and three,
so the polynomial is not a local square. Hence $c$ is integral at 3.
The polynomial $-720c^2+2152c+275$ has roots of valuations $-2,0$
in $\Q_3$. At the first root the ratio of differences lies
in $1+9\Z_3$ and is square. Thus $\lambda$ is a
$\Q_3$-square.

At 5 a negative even $v_5(c)$ would make $720c^2$ the unique term
of lowest, odd valuation in $q_T(c)$. Hence $c$ is integral there.
The same quadratic has $\Q_5$ roots of valuations $-1,2$.
The ratio at the first root lies in $1+5\Z_5$, so
$\lambda$ is a $\Q_5$-square. Since $q_A(c)>0$ and
$q_A(c_0)>0$, both parameters lie strictly between its two real
roots. Their root ratios are positive at both embeddings, forcing
$\lambda>0$. The known valuations leave only rational classes 1
and 2. The second is nonsquare in $\Q_3$. Thus $\lambda$
is a rational square and can be absorbed into $z$.

Let $s^2-s-1=0$, $r=(14+27s)/20$ and $\gamma=1/4-r$.
The normalized identity projected to $K=\Q(s)$ makes
$(c-r)/\gamma$ square, so
\[
 z^2=\gamma(c-r)(720c^2+120c+501).
\]
With $a=720,b=120,d=501,L=a\gamma$, its monic model is
\[
 E_0:v^2=u^3+\gamma(b-ar)u^2+L\gamma(d-br)u-L^2\gamma dr,
 \qquad(u,v)=(Lc,Lz),\quad\phi_0(u,v)=(u:L).
\]
The point $P_0=(L/4,24L\gamma)$ is non-torsion. The computation
transports both this point and the rational function through an exact
minimal-model isomorphism $\iota:E_0\to E$, so the function on $E$
is $\phi_0\circ\iota^{-1}$. It has the original $c$-coordinate as
its image. The rank over $K$ is one and torsion has order two.
Two-saturation and torsion give finite odd index; elliptic Chabauty
returns three elements, index obstruction 8 and images
$\{1/4,\infty\}$. The coprime-index condition proves completeness.
At $1/4$ the two paired-center squares are 1 and 4, whose chains
repeat $\pm1/2,\pm3/2$, including the central chain. Infinity is
not a center parameter. All auxiliary class-group proof bounds are
Minkowski bounds. The certificates are
\cert{check_type4_exclusion.m} and
\cert{check_type4_scalar_normalization.py}.

\subsection{The complete elliptic quotient for Type 2}\label{detail:type2}
The necessary Type~2 relation for the rational tripleton centers $u,v$ is
\[
 8u^4-40u^2v^2-58u^2+80v^4+50v^2+95=0.
\]
Consequently $y=8u^2-20v^2-29$ gives a point on
\[
 \mathcal D:y^2=-240v^4+760v^2+81.
\]
The smooth projective model has no rational infinities since $-240$
is nonsquare. For $v\ne0$, the maps
\[
 X=\frac{9(y+9)}{8v^2}+\frac{95}{2},\quad Y=\frac{9X}{2v},
 \qquad v=\frac{9X}{2Y},\quad
 y=\frac89(X-95/2)v^2-9
\]
identify it with $E:Y^2=X^3-95X^2+2560X$.
The omitted points $(0,\pm9)$ map to $O,(0,0)$ respectively.
The maps extend to the smooth projective models.
The dual curve is $E':Y^2=X^3+190X^2-1215X$.

Use \eqref{eq:early-isogeny-quartic} and
\eqref{eq:early-isogeny-rank}. On $E$ the quadratic
$X^2-95X+2560$ is positive on the real line, since its discriminant
is $-1215$. Thus any nonzero real abscissa is positive and the
candidate classes are $1,2,5,10$. The classes 1 and 10 occur at
the two isogeny-kernel points. On $E'$ the classes are
$\pm1,\pm3,\pm5,\pm15$, with 1 and $-15$ occurring at its
kernel points. Every other class is excluded by this complete table
of primitive-chart right-side residues.
\begin{center}
\begin{tabular}{ccll}
\toprule
Curve & $d$ & Modulus & Residues\\
\midrule
$E$ & $2,5$ & $3$ & $2$\\
$E'$ & $-1,15$ & $16$ & $7,12,15$\\
$E'$ & $3,-5$ & $16$ & $3,11,12$\\
$E'$ & $-3$ & $128$ & $5,21,29,69,80,93,117,125$\\
$E'$ & $5$ & $128$ & $5,13,21,77,80,85,101,109$\\
\bottomrule
\end{tabular}
\end{center}
For each modulus $p^e$, the charts $(a,1)$ and $(1,b)$, $p\mid b$,
exhaust all primitive pairs up to a unit. Every listed residue is
nonsquare. Both descent images have order two, and the rank formula
therefore proves rank zero.

The integral discriminant is $-2^{22}3^5 5^3$. Good reduction gives
$\#E(\F_7)=\#E(\F_{11})=10$, bounding all rational
torsion by ten, including the two possibly exceptional primary parts.
The ten rational points
\[
 O,(0,0),(16,\pm144),(40,\pm120),(64,\pm192),(160,\pm1440)
\]
attain this bound and exhaust $E(\Q)$. The inverse formulas
give exactly
\[
 (v,y)=(0,\pm9),\quad(\pm1/2,\pm16),\quad(\pm3/2,\pm24),
\]
with the two signs independent. At $v=0$, the recovered value
$u^2=(y+20v^2+29)/8$ is $19/4$ or $5/2$, both nonsquares.
At $v=\pm1/2,\pm3/2$, a $B$ tripleton contains a central $A$
vertex $\pm1/2$. Thus no point of the complete quotient lifts.
The independent checker \cert{check_type2_exclusion.py} verifies the
maps, complete residue tables, point counts, point list and substitutions.
The additional check \cert{check_type2_exclusion.m} is independent of
the elementary rank proof.

\subsection{The three seven-adic disks for Type 14}\label{detail:type14}
For Type~14 the moment branch is $a=(5b+3)/3$, where $a$ is the
$A$ doubleton center square and $b$ the $B$ doubleton center square.
The other three $B$ squared centers have power sums
$p_k=C_3(k)+M_2(k,a)-M_2(k,b)$, where
$M_L(k,z)=\sum_{d\in O_L}(\sqrt z+d)^{2k}$ and
$C_3(k)=\frac12\sum_{d\in O_6}d^{2k}$.
Odd powers of $\sqrt z$ cancel, so these are rational polynomials.
The monic singleton cubic is determined by Newton's formulas
$e_1=p_1$, $e_2=(p_1^2-p_2)/2$,
$e_3=(p_1^3-3p_1p_2+2p_3)/6$.
Its discriminant is
$-4bh_3(b)/243$, where
\[
 h_3(b)=900b^5-7100b^4-21675b^3+71658b^2+3321b-93312.
\]
Since $b=v^2\ne0$ and the singleton cubic splits rationally, a valid
input gives a finite point on the smooth genus-two curve
\[
 C:y^2=-3h_3(b).
\]
The points $\infty,(9,\pm3240)$ are rational. The Jacobian
has rank at most one by the unconditional two-descent in
\cert{check_type14_exclusion.m}. The class-group proof callback is
the actual Minkowski bound for every encountered order. In particular
the quintic root field has discriminant $-7745836638000$ and proof
bound 136074. We establish rank one and the complete list of eligible
points by one local calculation.

Put $a_0=-2700$, $X=a_0b$, $Y=a_0^2y$ and
$F(X)=a_0^4(-3h_3(X/a_0))$. The monic model $Y^2=F(X)$ has smooth
reduction at seven, with only $\infty,(4,\pm3)$ as residue points.
The three rational points above already occupy those three disks.
Let $D=[(9,3240)-\infty]$, viewed on the monic model. Exact Cantor
arithmetic for the reduced Mumford representative $(U,V)$ of $25D$
gives
\[
\begin{gathered}
 \deg U=2,\quad U\equiv(X-5)^2\pmod7,\quad
 v_7(\operatorname{disc}U)=2,\quad
 \operatorname{disc}(U)/7^2\equiv5\pmod7,\\
 7V\equiv5X+3\pmod7,\qquad F(5)\equiv5\pmod7.
\end{gathered}
\]
These congruences are independently recomputed by
\cert{check_type14_seven_adic.py}. They include sufficient precision
to determine the logarithm direction.

Because the discriminant unit 5 is nonsquare, the two support
abscissas $x_1,x_2$ are in the unramified quadratic extension of
$\Q_7$, have difference of valuation one and both reduce
to 5. The ordinates $y_i=V(x_i)$ are units. The leading coefficient
of $V$ has valuation $-1$, so $y_1-y_2$ is a unit; its two reductions
are the opposite roots of 5. Hence $P_1=(x_1,y_1)$ and
$\iota P_2=(x_2,-y_2)$ lie in the same disk, and
$25D=[P_1-\iota P_2]$. For
$\omega_0=dX/(2Y)$ and $\omega_1=X\,dX/(2Y)$, the local integrals are
\[
 \left(\int_{\iota P_2}^{P_1}\omega_0,
       \int_{\iota P_2}^{P_1}\omega_1\right)
 \equiv\frac{x_1-x_2}{-2y_2}(1,x_2)\pmod{7^2}.
\]
Indeed $Y$ is a unit and its expansion in $X-x_2$ is integral.
For every primitive term of degree $k\geq2$, the valuation is at
least $k-v_7(k)\geq2$. The first logarithm coordinate has valuation
one and is nonzero, proving $D$ non-torsion and the rank exactly one.
If $\lambda$ is the second coordinate divided by the first, then
$\lambda\in\Z_7$ and $\lambda\equiv5\pmod7$.
The differential $\omega_1-\lambda\omega_0$ therefore annihilates
the entire rational Mordell--Weil group.

At each finite residue point $X=4\ne5$, this differential has unit
coefficient in the local coordinate $X-4$. At infinity $\omega_1$
is a unit and $\omega_0$ vanishes. Thus in all three rational residue
disks it has integral coefficients and unit constant coefficient.
Its primitive is injective on $7\Z_7$. To see this for any
two distinct parameters $t_1,t_2$, the linear term in their primitive
difference has valuation $v_7(t_1-t_2)$, whereas each degree-$k$
term, $k\geq2$, has valuation at least
\[
 v_7(t_1-t_2)+k-1-v_7(k)>v_7(t_1-t_2).
\]
For two rational points in a disk the primitive difference must vanish
because their Abel--Jacobi difference is a rational Jacobian point.
Each disk therefore contains at most its known rational anchor.
This proves $C(\Q)=\{\infty,(9,\pm3240)\}$.
The rank-one Chabauty intrinsic in the Magma certificate independently
returns the same set.

For a genuine input $b$ is finite, so $b=9$ and $a=16$. The
singleton polynomial becomes
\[
 K(X)=\frac{(4X-1)(4X-9)(4X-81)}{64}.
\]
The $A$ doubleton center has absolute value 4, the $B$ doubleton
center has absolute value 3, and its singleton absolute centers are
$1/2,3/2,9/2$. With the prescribed offsets both inputs are exactly
\[
 \{\pm1/2,\pm3/2,\pm5/2,\pm7/2,\pm9/2\},
\]
contradicting disjointness. The exclusion covers every rational height.
